\section{The compact wall correspondences}\label{sec:walls}

The central factor classification determines the partial wall
targets as well as the central target. We first identify these compact
maps, then compute their actions. This is where the global moduli
problem enters the wall calculation.

\subsection{Partial targets and whole fibres}
Choose a small rational nonzero parameter $\eta_*=(s_*,t_*)$ on
one wall $\zeta_i=0$. Choose generic polarizations in the adjacent
sectors, close enough that there is no intervening wall. The finite
wall theorem and the exhaustive nearby list imply that every bundle
in either chamber is $\eta_*$-semistable. Normalize the common
image of their projective slope classifying maps. As in the central
construction, connected fibres and normal-function descent give a
normal projective partial target $Z_i$ and morphisms
\begin{equation}\label{eq:partial-target}
 M_+\longrightarrow Z_i\longleftarrow M_-.
\end{equation}
Their fibres are the $\eta_*$-graded classes. This uses the same
slope classifying construction of Greb--Toma \cite{gt} and the
whole-fibre descent in \cite[\S4.5]{wuebben17}.

There are no additional equal-slope decompositions on this partial
wall. Indeed the $H$-Jordan--H\"older classes of a proper equal-slope
factor must be a subcollection of $a_0,a_1,a_2$, or a varying
$w$-factor. Its equality equation is the corresponding subset sum
of the $\zeta_i$. At a nonzero point just one of these three
equations vanishes. The bounded wall theorem excludes further walls
after shrinking the neighborhood of $H$. Local freeness and the
exclusion of lost charge persist on the partial wall by the
central factor and charge result recalled above.

\begin{proposition}\label{prop:compact-walls}
If $i=0$ or $2$, the two morphisms in \eqref{eq:partial-target}
are the two resolutions of one isolated elementary Mukai contraction;
their birational map is an elementary Mukai flop. If $i=1$, the
reduced singular locus of $Z_1$ is a smooth compact K3 surface, and
both morphisms are its intrinsic blow-up. In particular they define
a unique isomorphism $M_+\simeq M_-$ over $Z_1$.
\end{proposition}

\begin{proof}
For $i=0,2$, the only strictly polystable class is $L_i\oplus G$,
where $G$ has vector $v-a_i$. The other two line factors have unequal
$\eta_*$-slopes. Their one-dimensional cross Ext gives the unique
nontrivial extension in the stable direction. It is stable: a
destabilizing subsheaf would, by the exhaustive nearby wall list,
give the wrong ordered line subfactor. Its vector is spherical.
Uniqueness also follows from stability and Riemann--Roch: two
nonisomorphic stable sheaves of the same spherical vector have
zero Hom in both directions and hence nonpositive Euler
characteristic, whereas that characteristic is two.

The factors $L_i,G$ have equal slope and are stable of different
ranks. Thus cross Hom and second Ext vanish, and
\[
 \dim\Ext^1(L_i,G)=(a_i,v-a_i)=3.
\]
They are rigid, so the full normal quotient is the doubled
three-arrow quotient without a tangent factor. Its two resolutions
are the incidence resolutions of the cone of rank-one square-zero
endomorphisms of $\C^3$.

These are whole compact fibres. To see why the local assertion
suffices here, take a contrary sequence of points in a wall fibre
outside its fixed classifying chart as the transverse parameter
tends to zero. No-charge compactness gives a smooth gauge limit.
Normalize the extension inclusion or quotient in $L^2$ and pass to
its nonzero limiting map. Equal-slope polystability splits off the
corresponding stable factor, and the determinant and factor list
identify the other. The limit is $L_i\oplus G$, contradicting
exclusion from the chart. The finite zero-section core has radius
of order the square root of the transverse moment parameter, so
the whole finite plane is in this chart as well. The two
classifying maps are isomorphisms off that fibre. The common
blow-up is locally the total space of $\OO(-1,-1)$ over the
incidence flag variety in $\PP^2\times(\PP^2)^*$. Its flag
variety is the exceptional divisor of the four-dimensional
blow-up. Gluing to the common regular open set proves the
compact elementary-flop assertion.

For $i=1$, write $A_1=L_1$. The polystable objects are $A_1\oplus F$
with $F$ stable of vector $w$ at $\eta_*$. Since
$\zeta_0=-\zeta_2\ne0$, this is a generic polarization for the
auxiliary problem. Its full moduli is the smooth compact K3 surface
$S$, including the curve over the central split class. Equal slopes
and different ranks imply cross Hom vanishing at every point of
$S$. Serre duality therefore gives constant cross-extension
dimension $(a_1,w)=2$.

The full classifying model along this locus is the product of a
smooth surface germ with $\C^2/\{\pm1\}$. The effective scalar
stabilizer acts trivially on the tangent directions to $S$. The
same compactness and normalized-map argument covers the entire
$\PP^1$ fibre over each point of $S$; compactness of $S$ gives a
finite cover of these models. The reduced singular locus is $S$,
and there is no deeper singularity on this nonzero wall.

The minimal resolution of the $A_1$ surface is the blow-up of
its reduced singular point. Blow-up commutes with the smooth
local product, so both maps are locally $\Bl_S Z_1\to Z_1$.
Their identifications agree on the dense regular set and hence
on overlaps. They glue to a unique global identification with
the intrinsic blow-up. Proper GAGA makes this identification
algebraic. Local universal $w$-families may be twisted, but all
these constructions use their coarse classifying spaces and
descend; no splitting of the algebraic Brauer class is required.
The resulting isomorphism also commutes with the maps to $Z$,
since its two composites agree on the dense stable open and $Z$
is separated.
\end{proof}

\subsection{The divisor sign and the reflection}
Use $+$ for the side $\zeta_1<0$ and $-$ for $\zeta_1>0$ in this
subsection. Let $D_\pm$ be the exceptional divisor of the maps
to $Z_1$, and let $\theta_\pm$ denote the Mukai markings.

\begin{lemma}\label{lem:opposite-divisors}
The two effective divisor classes are
\[
 [D_+]=\theta_+(a_1),\qquad [D_-]=-\theta_-(a_1).
\]
Both determinant sections have multiplicity one.
\end{lemma}

\begin{proof}
The first formula is the determinant calculation of
\cite[\S7.2]{wuebben17}, applied over the full auxiliary $S$ at
this partial wall. We give the opposite-side argument explicitly.
On the minus side $\mu(A_1)<\mu(E)$, so $\Hom(E,A_1)=0$.
The locus $\Hom(A_1,E)\ne0$ is exactly $D_-$.
For a nonzero map, let $B$ be the saturated rank-one hull
of its image. The quotient $E/B$ is torsion-free, so the
depth lemma on the smooth surface makes $B$ reflexive,
hence a line bundle. The effective zero divisor of
$A_1\to B$ has degree zero by $H$-semistability of $E$;
ampleness of $H$ forces it to vanish. Thus $B=A_1$.
The torsion-free quotient $F$ is $H$-semistable of vector
$w$, and \cite[Lemma~4.1]{wuebben17} makes it locally free.
It is semistable at the partial wall because $E$ is and
$A_1$ has the same wall slope; the auxiliary polarization
is generic, so $F$ is stable there. Conversely every
nontrivial extension on this side has an inclusion of $A_1$.

The perfect complex $Rp_*R\Hom(A_1,\cE)$ has amplitude
$[0,1]$ and virtual rank zero, since $(a_1,v)=0$. Its
determinant section is a section of the inverse determinant
line. The first Chern class of that line is
$-\theta_-(a_1)$. At an extension
$0\to A_1\to E\to F\to0$ the kernel has dimension one:
$\Hom(A_1,E)=\C$, and hence $\dim\Ext^1(A_1,E)=1$.
The derivative of the determinant is the obstruction to
preserving the inclusion. The map
\begin{equation}\label{eq:det-transverse}
 \Ext^1(E,E)\longrightarrow\Ext^1(A_1,E)
\end{equation}
is surjective. In the long exact sequence its cokernel maps
into $\Ext^2(F,E)\to\Ext^2(E,E)$, which is an isomorphism:
the Serre-dual map $\Hom(E,E)\to\Hom(E,F)$ is generated by
the quotient and is an isomorphism. Thus a local two-term
presentation of the index has corank one with a nonzero
normal derivative of its determinant. The zero is simple.
This proves the second formula with multiplicity one; the
dual inclusion/quotient argument gives the same multiplicity
on the plus side.
\end{proof}

\begin{proposition}\label{prop:marked-wall}
In the Mukai markings the isomorphism across $\zeta_1=0$
acts on degree two by $s_{a_1}$ and on the integral ring
by Markman's $\gamma_{s_{a_1}}$. Across $\zeta_0=0$ or
$\zeta_2=0$, the corrected flop correspondence acts by
$\gamma_{\id}$.
\end{proposition}

\begin{proof}
Let $f:M_+\to M_-$ be the divisorial isomorphism and set
$\psi=\theta_-^{-1}f_*\theta_+$. The universal families
agree up to a base line on the common stable open. Its
complement is the connected smooth prime divisor $D_-$,
so the kernel of restriction in degree two is $\Z[D_-]$,
by the Thom sequence. Lemma~\ref{lem:opposite-divisors}
implies $\psi(a_1)=-a_1$, while
$\psi(x)-x$ is an integral multiple of $a_1$ for every
$x\in v^\perp$. Write it as $k a_1$. Preservation of the
Beauville--Bogomolov form gives
\[
 (x,a_1)=(\psi x,\psi a_1)=-(x,a_1)+2k,
\]
so $k=(x,a_1)$, as asserted.

For an isolated plane wall let $\Gamma$ be the graph
closure. Rie\ss\ \cite[Theorem~3.2, Lemma~3.5 and \S6]{riess}
gives the integral multiplicative correspondence
\begin{equation}\label{eq:flop-correspondence}
 \mathcal Z=[\Gamma]+[P\times P'],\qquad
 \mathcal Z_*z=\Gamma_*z+\left(\int_Pz\right)[P']
 \quad(z\in H^4).
\end{equation}
Its transpose is the inverse. On the complement of the
plane the universal families agree up to a base line.
Deleting complex codimension two does not change degree-two
cohomology, so $\mathcal Z_*$ is the identity in the Mukai
markings. The rational cohomology of a K3$^{[2]}$ fourfold
is generated by degree two, as is also immediate from the
formulas in Section~\ref{sec:ring}. Thus both ring maps
agree rationally with the indicated Markman maps. Their
integral cohomology is torsion-free, so the agreement is
integral. This reasoning uses the corrected correspondence,
not the graph alone.
\end{proof}

\subsection{The six effective configurations}
The cyclic order of the chambers and their wall maps is shown in
Figure~\ref{fig:chambers}. The complete effective cycle table is
Table~\ref{tab:cycles}. The notation $p_i$ here means the class
obtained from the plane formula \eqref{eq:plane-formula};
Section~\ref{sec:ring} identifies it with
\eqref{eq:planes-intro} in the common ring $R$.

\begin{figure}[htbp]
\centering
\begin{adjustbox}{max width=\textwidth,center}
\begin{tikzpicture}[>=Stealth,
 chamber/.style={draw,rounded corners=2pt,fill=blue!3,align=center,
 minimum width=21mm,minimum height=10mm,font=\small},
 wall/.style={font=\footnotesize,fill=white,inner sep=2pt}]
\node[chamber] (A) at (0,2.1) {$A$\\$(+,-,+)$};
\node[chamber] (B) at (3.8,1.05) {$B$\\$(+,-,-)$};
\node[chamber] (C) at (3.8,-1.05) {$C$\\$(+,+,-)$};
\node[chamber] (D) at (0,-2.1) {$D$\\$(-,+,-)$};
\node[chamber] (E) at (-3.8,-1.05) {$E$\\$(-,+,+)$};
\node[chamber] (F) at (-3.8,1.05) {$F$\\$(-,-,+)$};
\draw[->,thick] (A) -- node[wall,above,sloped] {$\zeta_2=0,\ \id$} (B);
\draw[->,thick,dashed] (B) -- node[wall,right] {$\zeta_1=0,\ s_{a_1}$} (C);
\draw[->,thick] (C) -- node[wall,below,sloped] {$\zeta_0=0,\ \id$} (D);
\draw[->,thick] (D) -- node[wall,below,sloped] {$\zeta_2=0,\ \id$} (E);
\draw[->,thick,dashed] (E) -- node[wall,left] {$\zeta_1=0,\ s_{a_1}$} (F);
\draw[->,thick] (F) -- node[wall,above,sloped] {$\zeta_0=0,\ \id$} (A);
\node[align=center,font=\small] at (0,0) {four plane flops\\two divisorial isomorphisms};
\end{tikzpicture}

\end{adjustbox}
\caption{The chamber adjacency cycle. Solid edges are elementary
Mukai flops with their corrected cohomology correspondences; dashed
edges are the divisorial isomorphisms. The edge labels give the
vanishing coordinate and the Mukai action.}
\label{fig:chambers}
\end{figure}

\begin{table}[htbp]
\caption{Effective deepest surfaces and curves. In a ruled row the
last three entries are the plane line, the common negative section,
and the ruling. Here $f=c_0+c_2$.}
\label{tab:cycles}
\centering
\begin{adjustbox}{max width=\textwidth,center}
\begin{tabular}{cccl}
\toprule
Chamber & Signs $(\zeta_0,\zeta_1,\zeta_2)$ & Surfaces & Curves\\
\midrule
$A$ & $(+,-,+)$ & $p_0,p_2$ & $c_0,c_2$ (two planes)\\
$B$ & $(+,-,-)$ & $p_2,p_0-p_2$ & $-c_2;\ -c_2;\ f$\\
$C$ & $(+,+,-)$ & $p_0,p_2-p_0$ & $c_0;\ c_0;\ -f$\\
$D$ & $(-,+,-)$ & $p_0,p_2$ & $-c_0,-c_2$ (two planes)\\
$E$ & $(-,+,+)$ & $p_2,p_0-p_2$ & $c_2;\ c_2;\ -f$\\
$F$ & $(-,-,+)$ & $p_0,p_2-p_0$ & $-c_0;\ -c_0;\ f$\\
\bottomrule
\end{tabular}
\end{adjustbox}
\end{table}

Here is the geometric calculation of the table. If $P_2$ is
flopped, its graph image is $-2[P_2']$ and the exceptional
product contributes $3[P_2']$; see
\cite[Example~6.6]{llw} or \cite[\S6]{riess}. Thus the
corrected image is $[P_2']$. The other plane meets $P_2$
transversely once. Its graph image is its blow-up at that
point, and the product term adds $[P_2']$. Consequently
\begin{equation}\label{eq:surface-transform}
 p_2\longmapsto[P_2'],\qquad
 p_0\longmapsto[\Sigma]+[P_2'],\qquad \Sigma\simeq\FF_1.
\end{equation}
The common plane line and negative section represent $-c_2$.
A line in $P_0$ missing the blown-up point keeps its class;
on $\FF_1$ it is the sum of negative section and ruling.
The ruling therefore has class $f=c_0+c_2$. This gives row
$B$ and the intersection matrix
\begin{equation}\label{eq:ruled-matrix}
 \begin{pmatrix}p_2^2&p_2(p_0-p_2)\\
 p_2(p_0-p_2)&(p_0-p_2)^2\end{pmatrix}
 =\begin{pmatrix}3&-2\\-2&4\end{pmatrix}.
\end{equation}
The reflection sends $c_0$ to $-c_2$, $c_2$ to $-c_0$,
and $f$ to $-f$; it exchanges $p_0,p_2$ and preserves the
tangent Chern class. This proves row $C$. The next flop
gives the two-plane row $D$, with both line signs reversed.
Repeating gives $E,F$ and returns to $A$.

\begin{figure}[htbp]
\centering
\begin{adjustbox}{max width=\textwidth,center}
\begin{tikzpicture}[>=Stealth,font=\small]
\draw[fill=blue!7,thick] (-4,0) ellipse (1.15 and .7);
\draw[fill=blue!7,thick] (-1.7,0) ellipse (1.15 and .7);
\fill (-2.85,0) circle (2pt);
\node at (-4,0.15) {$P_0$};\node at (-1.7,0.15) {$P_2$};
\node at (-4,-.3) {$p_0,\ c_0$};\node at (-1.7,-.3) {$p_2,\ c_2$};
\node at (-2.85,1.1) {chamber $A$};
\draw[->,thick] (-.35,0) -- node[above] {flop $P_2$} (1.45,0);
\draw[fill=blue!7,thick] (1.8,-.65) rectangle (3.3,.65);
\draw[fill=orange!8,thick] (3.3,-.65) rectangle (5.5,.65);
\draw[very thick,blue] (3.3,-.65) -- (3.3,.65);
\node at (2.55,.25) {$P_2'$};\node at (2.55,-.25) {$p_2$};
\node at (4.4,.25) {$\Sigma\simeq\FF_1$};\node at (4.4,-.25) {$p_0-p_2$};
\node at (3.65,1.1) {chamber $B$};
\node[align=center] at (3.3,-1.05) {common curve $-c_2$};
\node[align=center] at (-2.85,-1.5) {moving fibre $f\rightsquigarrow c_0+c_2$};
\node[align=center] at (3.65,-1.65) {ruling $f=c_0+c_2$};
\end{tikzpicture}

\end{adjustbox}
\caption{The two configurations and their classes across the flop
from $A$ to $B$. The right-hand common curve is a line in $P_2'$
and the negative section of $\Sigma$. The moving fibre specializes
in $A$ to the reduced sum $c_0+c_2$; in $B$ this class is the
ruling of $\Sigma$. The drawings indicate incidences and cycle
classes, not metric shapes.}
\label{fig:configurations}
\end{figure}

\begin{corollary}\label{cor:gallery}
The chosen cohomological correspondences along any two paths in
the six-chamber adjacency graph with the same endpoints agree.
The complements $V_\sigma$ are isomorphic over $Z\setminus\{o\}$,
and the wall actions induce their cohomology identifications.
\end{corollary}

\begin{proof}
The six forward actions in the Mukai marking are
$\id,s_{a_1},\id,\id,s_{a_1},\id$. Their product is the identity.
Each reverse edge is the inverse correspondence. Cancellation
of reverse edges and insertion of complete six-cycles generate
the relations between paths in this graph. The ring assertion
follows. For the complements, the central map is an isomorphism
over $U$ and has the unique relative minimal $A_1$ resolution
along $S^\circ$. The local identifications agree on the dense
regular locus and glue over $Z\setminus\{o\}$.
\end{proof}

The corollary concerns these ordinary cohomological actions. It
does not identify compositions of Chow correspondences or derived
functors, and it is not a monodromy calculation for a smoothing.
